\chapter{Нейрони за кількістю дендритів}
\chaptermark{Кількість дендритів}
\label{sec:dendrites}
\begin{chaptergoals}
\item чому ємність мембрани робить нейрон із небагатьма входами швидким, а з багатьма --- повільним;
\item як датчики, буфери й детектори з двома дендритами тримають окремий сигнал швидким і точним;
\item чому змішувачі з близько чотирма входами дають великому суматору зберегти багато відповідностей;
\item як великі дерева додають і класифікують, а дендрити-виходи роблять із клітини кілька детекторів.
\end{chaptergoals}

\section{Кількість входів як параметр схеми}

У розділі~\ref{sec:neurontypes} ми сортували нейрони за тим, що викидає їхній вихід, а в розділі~\ref{sec:eetypes} --- за тим, яким приладом вони працюють. Є й третя ознака, яку інженер помітить одразу: \emph{скільки в нейрона входів}. Входи збирають дендрити --- гілки, на яких сідають чужі аксони (розділ~\ref{sec:dofa}, підрозділ~\ref{sec:weightstore}). Одні нейрони не мають дендритів зовсім, інші мають один чи два, ще інші --- дерево, що збирає сотню тисяч входів. Для електронника це навантажувальна здатність за входом, fan-in, і вона майже завжди відповідає задачі.

Чому кількість і розмір дендритів такі важливі, видно з простої оцінки. Мембрана --- конденсатор із питомою ємністю $c_m\approx1$~мкФ/см$^2$, і що більша площа $A$ нейрона разом із його дендритами, то більший заряд треба принести, щоб підняти напругу на $\Delta V$. Час, за який вхідний струм $I$ це зробить, якщо не враховувати витоку,
\begin{empheq}[box=\keybox]{equation}
t \;\approx\; \frac{c_m\,A\,\Delta V}{I} .
\label{eq:dendriteload}
\end{empheq}
Маленький нейрон із кількома короткими дендритами має площу порядку $2\cdot10^{-5}$~см$^2$ і ємність близько $20$~пФ. Один великий синапс зі струмом у кілька наноамперів піднімає його на $20\mV$ менш ніж за десяту частку мілісекунди. У нейрона з великим деревом площа разів у п'ятнадцять більша, а ємність близько $300$~пФ; звичайний синапс зі струмом близько $0.1$~нА зарядив би його за $60\ms$ --- набагато довше за сталу часу витоку, тобто не зарядив би ніколи. Такому нейронові потрібні десятки одночасних входів. Числа тут --- порядки величин, але висновок від них не залежить: \emph{мало входів --- швидко й точно, багато входів --- повільно й за сумою}.

\begin{figure}[t]
\centering
\begin{tikzpicture}[>=Stealth,
  soma/.style={circle,draw,very thick,fill=blue!6,minimum size=0.55cm,inner sep=0pt},
  ax/.style={->,very thick,red!70!black},
  den/.style={very thick,blue!60!black}]
% 0 dendrites: sensor on a line
\begin{scope}[xshift=0cm]
\draw[den,decorate,decoration={snake,amplitude=1pt,segment length=4pt}] (-0.9,0) -- (0,0);
\node[font=\scriptsize] at (-0.9,0.3) {датчик};
\draw[ax] (0,0) -- (1.0,0);
\draw[thick] (0.1,0) -- (0.1,0.45);
\node[soma,minimum size=0.4cm] at (0.1,0.65) {};
\node[font=\scriptsize,align=center] at (0.05,-0.75) {$0$\\лінія з датчиком};
\end{scope}
% 1 dendrite
\begin{scope}[xshift=3.3cm]
\draw[den] (-0.9,0) -- (-0.28,0);
\node[soma] at (0,0) {};
\draw[ax] (0.28,0) -- (0.9,0);
\node[font=\scriptsize,align=center] at (0,-0.75) {$1$\\буфер, повторювач};
\end{scope}
% 2 dendrites: one per ear
\begin{scope}[xshift=6.2cm]
\draw[den] (-0.28,0) -- (-0.9,0); \draw[den] (0.28,0) -- (0.9,0);
\node[soma] at (0,0) {};
\draw[ax] (0,-0.28) -- (0,-0.6);
\node[font=\scriptsize] at (-0.9,0.25) {ліве}; \node[font=\scriptsize] at (0.9,0.25) {праве};
\node[font=\scriptsize,align=center] at (0,-1.0) {$2$\\«І» в часі};
\end{scope}
% ~4 short dendrites: mixer
\begin{scope}[xshift=9.0cm]
\foreach \a in {45,135,225,315}{\draw[den] (\a:0.28) -- (\a:0.7); \fill[blue!60!black] (\a:0.72) circle (0.06);}
\node[soma] at (0,0) {};
\draw[ax] (0.28,0) -- (1.0,0);
\node[font=\scriptsize,align=center] at (0,-1.0) {$\sim4$\\змішувач};
\end{scope}
% many: summing tree
\begin{scope}[xshift=12.0cm]
\foreach \a in {60,90,120}{\draw[den] (0,0.28) -- ++(\a:0.55) coordinate (b); \foreach \d in {-20,20}{\draw[den] (b) -- ++(\a+\d:0.35);}}
\foreach \a in {-120,-60}{\draw[den] (0,-0.28) -- ++(\a:0.45);}
\node[soma] at (0,0) {};
\draw[ax] (0.28,0) -- (1.0,0);
\node[font=\scriptsize,align=center] at (0,-1.0) {$10^4$--$10^5$\\суматор};
\end{scope}
\end{tikzpicture}
\caption{П'ять класів за кількістю входів. Синім --- дендрити (входи), червоним --- аксон (вихід). Що менше входів, то швидше й точніше нейрон передає окремий сигнал; що більше, то краще він додає й класифікує. Це схема, а не малюнок форми клітин.}
\label{fig:dendriteclasses}
\end{figure}

\section{Нуль дендритів: датчик на кінці лінії}

Чутливі нейрони дотику, болю й розтягу м'язів (розділи~\ref{sec:sensors}, \ref{sec:pain}, \ref{sec:muscles}) не мають на тілі клітини жодного дендрита. Аксон виходить із тіла однією гілочкою і відразу ділиться надвоє: одна гілка йде до шкіри чи м'яза, і її кінець сам слугує датчиком; друга йде в спинний мозок. Імпульс біжить від датчика просто в спинний мозок, оминаючи тіло клітини. Для інженера це лінія зв'язку, у якої датчик стоїть на дальньому кінці, а тіло клітини висить збоку, як блок живлення: воно живить лінію, але в передаванні не бере участі. Вигода --- швидкість і надійність: на шляху немає жодного місця, де сигнал треба було б підсумовувати й знову вирішувати, чи передавати його.

Датчики світла й звуку --- ще крайніший випадок: це не нейрони-збирачі, а самі перетворювачі, у яких входом слугує не синапс, а білок-рецептор (рис.~\ref{fig:transducer}).

\section{Один дендрит: буфер}

Нейрон з одним дендритом і одним аксоном приймає одну вхідну лінію й передає її далі. Так улаштовані нюхові нейрони: єдиний дендрит виходить до поверхні носа й несе єдиний тип рецептора~\cite{Mombaerts2004}; так улаштовані передавальні клітини сітківки між датчиками світла й вихідними нейронами ока; так улаштовані нейрони, що зв'язують датчики вуха з мозком. Інженер назвав би їх буфером або повторювачем: вони не обчислюють, а доставляють один сигнал, іноді змінюючи його форму --- наприклад, перетворюючи плавну напругу датчика на імпульси, як у розділі~\ref{sec:sensors}.

\section{Два дендрити: «І» в часі}

Детектори збігів, що визначають напрямок на звук (рис.~\ref{fig:itd}), у ссавців часто мають рівно два дендрити, спрямовані в протилежні боки: на один надходять входи від лівого вуха, на другий --- від правого~\cite{Grothe2010}. Навіщо розносити входи? Згадаймо формулу~\eqref{eq:shunt}: коли на одній ділянці мембрани відкрито багато збуджувальних каналів, напруга там наближається до рівноважного потенціалу, і кожен наступний вхід додає дедалі менше. Тому багато входів від одного вуха, зібраних на одному дендриті, насичують його і не можуть самотужки довести нейрон до порога. А по одному входу з кожного дендрита додаються майже лінійно. Два дендрити перетворюють нейрон на суворіше «І»: спрацьовуй, лише якщо сигнали надійшли з обох боків і разом. На моделях показано, що таке рознесення справді поліпшує детекцію збігів~\cite{AgmonSnir1998}.

\section{Кілька коротких дендритів: змішувач і ретранслятор}

\emph{Змішувачі.} У колі навчання рухів із розділу~\ref{sec:motorlearning} близько сімдесяти мільярдів маленьких \nt{GLUT}-нейронів розширюють вхід в дуже довгий вектор. Кожен із них має лише близько чотирьох коротких дендритів, і кожен закінчується «пазуром», який охоплює закінчення одного входу. Кожен змішувач, отже, бачить лише близько чотирьох входів із багатьох і працює як маленький елемент «І» над своєю четвіркою. З $M$ вхідних ліній можна вибрати $\binom{M}{4}$ різних четвірок: за $M=100$ це майже чотири мільйони. Навіщо стільки? Пороговий елемент з $n$ входами може правильно поділити на два класи не більше ніж
\begin{empheq}[box=\keybox]{equation}
P_{\max} \;\approx\; 2n
\label{eq:cover}
\end{empheq}
випадкових образів~\cite{Cover1965}. Вихідний нейрон того самого кола отримує близько $10^5$ входів від змішувачів, і за~\eqref{eq:cover} верхня межа для нього --- порядку $2\cdot10^5$ різних відповідностей. Це межа для ідеального елемента: якщо всі ваги додатні, як у збуджувальних синапсів, і потрібен запас на шум, оцінка для цього самого нейрона --- близько $4\cdot10^4$ відповідностей~\cite{Brunel2004}. Змішувачі з малою кількістю входів потрібні саме для того, щоб великий суматор мав багато різних, майже незалежних входів~\cite{Marr1969,Albus1971}.

\emph{Ретранслятори.} На шляху від вуха до детекторів напрямку є нейрони з небагатьма короткими дендритами, які отримують лише кілька величезних синапсів, а деякі --- по суті один. За~\eqref{eq:dendriteload} це саме те, що треба: мала ємність і великий струм, і нейрон спрацьовує через частки мілісекунди після кожного вхідного імпульсу, майже не втрачаючи точності часу. Один із таких ретрансляторів отримує \nt{GLUT}, а видає \nt{GLYC}: це інвертор із точністю в десятки мікросекунд, що постачає швидке гальмування детекторам напрямку~\cite{BorstSoria2012}.

\section{Тисячі входів: суматор і класифікатор}

На другому кінці шкали --- \nt{GLUT}-нейрони кори з десятком тисяч входів і вихідні \nt{GABA}-нейрони кола навчання рухів із сотнею тисяч. Такий нейрон за~\eqref{eq:dendriteload} повільний і не може відгукнутися на окремий вхід: він додає. Це інтегрувальний нейрон розділу~\ref{sec:glutcomputer} і той суматор, з якого будують класифікатори. Велике дерево додає й нове: його гілки працюють як окремі підсхеми з власними порогами, тож один нейрон поводиться як маленька багатошарова мережа~\cite{Beniaguev2021,Gidon2020}.

\section{Дендрити, які є ще й виходом}

Є нейрони, що не мають аксона зовсім, а дендрити слугують і входом, і виходом: вони отримують сигнали й самі викидають медіатор. Найвивченіший приклад --- \nt{GABA}-нейрони сітківки, що беруть участь у визначенні напрямку руху (розділ~\ref{sec:gaba}, рис.~\ref{fig:direction}). У такого нейрона кожна гілка сама по собі відповідає на рух від тіла клітини до свого кінчика й видає гальмування з цього кінчика~\cite{Euler2002}. Один нейрон --- це кілька незалежних детекторів напрямку, по одному на гілку. Для інженера це не один елемент, а мікросхема з кількома каналами в одному корпусі.

\section{Підсумок}

\begin{table}[t]
\centering
\footnotesize
\setlength{\tabcolsep}{3pt}
\renewcommand{\arraystretch}{1.15}
\begin{tabular}{>{\raggedright}p{1.9cm}>{\raggedright}p{3.4cm}>{\raggedright}p{2.6cm}>{\raggedright}p{1.8cm}>{\raggedright\arraybackslash}p{3.8cm}}
\toprule
Дендритів & Приклад & Прилад & Входів & Що відомо \\
\midrule
$0$ & датчики дотику, болю, розтягу & лінія з датчиком на кінці & --- & встановлено \\
$1$ & нюхові нейрони, передавальні клітини сітківки, нейрони вуха & буфер, повторювач & $1$--небагато & встановлено \\
$2$ & детектори напрямку на звук & «І» в часі & два пучки & будову встановлено; виграш від рознесення входів показано на моделях \\
$\sim4$ & змішувачі кола навчання рухів & маленький елемент «І» & $\sim4$ & встановлено; роль розширення входу --- добре обґрунтована гіпотеза \\
небагато, короткі & ретранслятори на шляху від вуха & ретранслятор, інвертор & $1$--кілька величезних & встановлено \\
тисячі & \nt{GLUT}-нейрони кори, вихідні нейрони навчання рухів & суматор, класифікатор & $10^4$--$10^5$ & встановлено; гілки як підсхеми виміряно на окремих прикладах \\
дендрити-виходи & \nt{GABA}-нейрони напрямку в сітківці & багатоканальна мікросхема & багато & виміряно \\
\bottomrule
\end{tabular}
\caption{Нейрони за кількістю дендритів.}
\label{tab:dendrites}
\end{table}

\begin{keyblock}
\begin{proposition}[Кількість входів іде за задачею]
\label{prop:dendrites}
Де потрібні швидкість і точність окремого сигналу --- на датчиках, ретрансляторах і детекторах збігів, --- нейрони мають мало дендритів, мало входів і великі синапси: мала ємність~\eqref{eq:dendriteload} заряджається швидко. Де треба додавати й класифікувати, нейрони мають великі дерева з тисячами входів. Будову класів виміряно; обчислювальне пояснення добре обґрунтоване, але для багатьох класів лишається гіпотезою.
\end{proposition}
\end{keyblock}

Таблиця~\ref{tab:dendrites} доповнює два попередні сортування: за медіатором (таблиця~\ref{tab:types}) і за приладом (таблиця~\ref{tab:eetypes}). Усі три дивляться на той самий елемент із різних боків: що він видає, що він обчислює і скільки він слухає.

\section{Задачі}

\begin{exercise}[\lvlA]
\label{ex:19:1}
\emph{Скільки входів потрібно великому нейронові.} Нейрон із $C=300$~пФ і $\tau_m=20\ms$ отримує синапси --- джерела струму по $0.1$~нА, поріг на $20\mV$ вищий за спокій. (а)~Скільки синапсів має працювати одночасно, щоб дійти до порога взагалі; за $10\ms$; за $5\ms$? (б)~За який час дійде до порога нейрон із $C=20$~пФ від синапсу в $4$~нА?
\end{exercise}

\begin{exercise}[\lvlA]
\label{ex:19:2}
\emph{Чому чотири.} Змішувач --- «І» над $k$ лініями з $M=100$, в образі активна половина ліній, змішувачів $7\cdot10^{10}$. (а)~Скільки з них відповідає на образ за $k=2$; $4$; $40$? (б)~У скільки разів змішувачі збільшують кількість відповідностей~\eqref{eq:cover} вихідного нейрона з $10^5$ входами проти $M=100$ прямих ліній?
\end{exercise}

\begin{exercise}[\lvlB]
\label{ex:19:3}
\emph{Звідки береться $2n$.} Гіперплощина через початок координат ділить $P$ точок загального положення в $n$-вимірному просторі на два класи $C(P,n)=2\sum_{k=0}^{n-1}\binom{P-1}{k}$ способами. (а)~Доведіть, що за $P=2n$ роздільна рівно половина з $2^P$ розбиттів. (б)~Яка їхня частка роздільна за $n=100$ і $P=180$; $220$?
\end{exercise}

\begin{exercise}[\lvlB]
\label{ex:19:4}
\emph{Два дендрити як суворе «І».} Дендрит --- ділянка з витоком $g_L$, кожен вхід відкриває на ній провідність $g_0=0.5\,g_L$ з потенціалом $E_E$, тіло бачить $\kappa(V_1+V_2)$, поріг $\theta=1.1\,\kappa E_E$. Чому жодна кількість входів на одному дендриті не запускає нейрона і скільки входів потрібно на кожному?
\end{exercise}

\begin{exercise}[\lvlB]
\label{ex:19:5}
\emph{Проєктування ретранслятора.} Тремтіння імпульсу $\sigma_t=\sigma_V/\dot V$, де $\dot V$ --- нахил напруги біля порога. Потрібно $\sigma_t\le20$~мкс за шуму $\sigma_V=1\mV$; витоком знехтуйте. Скільки синхронних синапсів по $0.1$~нА треба для цього нейронові з $C=300$~пФ і яке тремтіння нейрона з $C=20$~пФ та одним синапсом у $4$~нА?
\end{exercise}

\begin{exercise}[\lvlB]
\label{ex:19:6}
\emph{Моделювання: заряд довгого дендрита.} Змоделюйте пасивний кабель $\tau_m\partial_tV=\lambda^2\partial_x^2V-V+i/g$ довжини $L=\ell/\lambda$ із закритими кінцями ($40$ ділянок, Ейлер). У дальній кінець подано короткий імпульс струму: коли й до якої частки напруги від того самого заряду, розмазаного по всьому дендриту, піднімається ближній кінець за $L=0.2$; $1$; $2$?
\end{exercise}

\begin{exercise}[\lvlC]
\label{ex:19:7}
\emph{Дослідження: оптимальна кількість входів.} Ємність $C=C_0+nc_s$, одночасно працюють $m$ входів зі струмом $I_s$, нейрон запам'ятовує $P\approx2n$ відповідностей~\eqref{eq:cover}. Як час відгуку залежить від $P$ за сталого $m$ і за сталої частки $m=\varphi n$? Запропонуйте критерій вартості й знайдіть оптимальне $n$ для ретранслятора й класифікатора.
\end{exercise}
